\section{Zadanie 3}
\autor{Dominik Korsa}
Ze standardowej talii 52 kart wyciągamy 13 (ręka brydżowa). Niech \(X\) to
liczba wyciągniętych asów, a \(Y\) to liczba wyciągniętych trefli.

\begin{enumerate}[label=(\roman*)]
  \item\label{set04.task03.1} Czy \(X\) i \(Y\) są niezależne?
  \item\label{set04.task03.2} Czy \(X\) i \(Y\) są nieskorelowane?
\end{enumerate}

\subsection*{Przypomnienie}
Zmienne losowe \(X\) i \(Y\) są \textbf{niezależne}, gdy dla każdego
\(x, y \in \mathbb{R}\):
\[
  P(X=x \land Y=y) = P(X=x) \cdot P(Y=y)
\]

Zmienne losowe \(X\) i \(Y\) są \textbf{nieskorelowane}, gdy:
\[
  E[X \cdot Y] = E[X] \cdot E[Y]
\]

\subsection*{Rozwiązanie podpunktu \ref{set04.task03.1}}
Nie są niezależne, pokażę kontrprzykład.

\begin{itemize}
  \item \(P(X=4)\) to prawdopodobieństwo, że wyciągniemy 4 asy, oczywiście
  \(P(X=4) > 0\).

  \item \(P(Y=13)\) to prawdopodobieństwo, że wyciągniemy tylko trefle, podobnie
  \(P(Y=13) > 0\).

  \item \(P(X=4 \land Y=13)\) to prawdopodobieństwo, że wyciągniemy tylko trefle,
  w tym cztery asy. Takie zdarzenie jest niemożliwe, więc
  \(P(X=4 \land = Y=13) = 0\).
\end{itemize}
Czyli nie jest prawdą, że dla każdego \(x, y \in \mathbb{R}\) zachodzi
\[P(X=x \land Y=y) = P(X=x) \cdot P(Y=y),\]
więc zmienne losowe \(X\) i \(Y\) nie są niezależne.

\subsection*{Rozwiązanie podpunktu \ref{set04.task03.2}}
Niech $X_i$ to indykator zdarzenia, że $i$-ta wylosowana karta to as:
\[
  E[X_i=1] = \frac{4}{52} = \frac{1}{13} \quad\text{(dla \(i \in [13]\))}
\]
Z liniowości wartości oczekiwanej:
\[
  E[X]
  = E[X_1 + \cdots + E_{13}]
  = E[X_1] + \cdots + E[X_{13}]
  = 13 \cdot \frac{1}{13} = 1
\]

Analogicznie niech $Y_i$ to indykator zdarzenia, że $i$-ta karta to trefl:
\[
  E[Y] = \cdots = 13 \cdot \frac{1}{4} = \frac{13}{4}
\]

% Spróbujmy policzyć lub oszacować \(E[XY]\). Skorzystajmy z warunkowej wartości
% oczekiwanej:
% \begin{align*}
%   E[XY]
%   &= \sum_{x=0}^{4} E[XY | X=x] \cdot P(X=x)\\
%   &= \sum_{x=0}^{4} x \cdot E[Y | X=x] \cdot P(X=x)\\
% \end{align*}

% \(P(X=x) = \frac{\binom{4}{x} \cdot \binom{48}{13-x}}{\binom{52}{13}}\)
\(X_iY_i\) to indykator zdarzenia, że \(i\)-ta wylosowana karta to as trefl:
\(E[X_iY_i] = \frac{1}{52}\).

Dla \(i \neq j\) \(X_iY_j\) to indykator zdarzenia, że \(i\)-ta wylosowana karta
jest asem, a \(j\)-ta to trefl. Te zdarzenia są niezależne, więc:
\[
  E[X_iY_j] = P(X_i = 1 \land Y_j = 1)
  = P(X_i = 1) \cdot P(Y_j = 1)
  = \frac{1}{4} \cdot \frac{1}{13} = \frac{1}{52}
\]

Czyli dla każdej pary \((i, j) \in [13] \times [13]\) zachodzi \(E[X_iY_i] = \frac{1}{52}\).

Z liniowości wartości oczekiwanej:
\[
  E[XY]
  = E\left[ \sum_{(i,j) \in [13] \times [13]} X_iY_j \right]
  = \sum_{(i,j) \in [13] \times [13]} E[X_iY_j]
  = 13^2 \cdot \frac{1}{52}
  = \frac{13}{4},
\]
Wiemy też, że:
\[
  E[X] \cdot E[Y] = 1 \cdot \frac{13}{4} = \frac{13}{4},
\]
więc warunek na \textbf{nieskorelowanie} zmiennych losowych \(X\) i \(Y\) zachodzi.

\subsection*{Odpowiedź}
\[
\boxed{\text{\(X\) i \(Y\) \textbf{nie są niezależne}, ale \textbf{są nieskorelowane}}}
\]
